% ============================================================
\chapter{Alg\`ebre Relationnelle}
\label{ch:algebre}
% ============================================================

\section{Introduction}

Quand vous \'ecrivez \texttt{SELECT nom FROM etudiants WHERE note > 15}, vous pensez en SQL. Mais derri\`ere cette syntaxe se cache un langage math\'ematique d'une pr\'ecision remarquable~: l'\emph{alg\`ebre relationnelle}, invent\'ee par Edgar F.\ Codd dans son article fondateur de 1970. Codd a eu l'intuition g\'eniale de voir les tables comme des \emph{relations} au sens math\'ematique (des sous-ensembles de produits cart\'esiens) et les requ\^etes comme des \emph{op\'erations} sur ces relations~: s\'election, projection, jointure, union, diff\'erence. L'alg\`ebre relationnelle est \`a SQL ce que la logique formelle est au langage naturel~: un socle rigoureux qui garantit que chaque requ\^ete a une s\'emantique pr\'ecise et non ambigu\"e.

Comprendre l'alg\`ebre relationnelle, c'est comprendre \emph{pourquoi} SQL fonctionne --- et comment l'optimiseur de requ\^etes transforme vos instructions en un plan d'ex\'ecution efficace.

\begin{definition}[Alg\`ebre relationnelle]
L'\emph{alg\`ebre relationnelle} est un ensemble d'op\'erateurs qui prennent
une ou deux relations en entr\'ee et produisent une relation en sortie.
C'est un langage \emph{proc\'edural}~: on sp\'ecifie \emph{comment}
obtenir le r\'esultat (s\'equence d'op\'erations).
\end{definition}

Les op\'erateurs se divisent en deux cat\'egories~:
\begin{itemize}
    \item \textbf{Op\'erateurs ensemblistes}~: union ($\cup$), intersection ($\cap$),
          diff\'erence ($-$), produit cart\'esien ($\times$).
    \item \textbf{Op\'erateurs sp\'ecifiques}~: s\'election ($\sigma$), projection ($\pi$),
          jointure ($\bowtie$), renommage ($\rho$), division ($\div$).
\end{itemize}

\section{S\'election ($\sigma$)}

\begin{definition}[S\'election]
La \emph{s\'election} $\sigma_\theta(R)$ retourne les tuples de $R$ qui satisfont
la condition $\theta$. La condition $\theta$ est une expression bool\'eenne
portant sur les attributs de $R$.
\[
\sigma_\theta(R) = \{ t \in R \mid \theta(t) = \text{vrai} \}
\]
\end{definition}

\begin{formules}[title=\textbf{Propri\'et\'es de la s\'election}]
\begin{align}
\sigma_{\theta_1 \land \theta_2}(R) &= \sigma_{\theta_1}(\sigma_{\theta_2}(R)) \tag{cascade}\\
\sigma_{\theta_1}(\sigma_{\theta_2}(R)) &= \sigma_{\theta_2}(\sigma_{\theta_1}(R)) \tag{commutativit\'e}\\
|\sigma_\theta(R)| &\leq |R| \tag{cardinalit\'e}
\end{align}
\end{formules}

\begin{exemple}
S\'electionner les \'etudiants de la fili\`ere Informatique~:
\[
\sigma_{\text{filiere} = \text{'Info'}}(\text{Etudiants})
\]

Cela correspond en SQL \`a~:
\begin{codeblock}[title=\textbf{\'Equivalent SQL de la s\'election}]
\begin{minted}{sql}
SELECT * FROM Etudiants WHERE filiere = 'Info';
\end{minted}
\end{codeblock}

\begin{codeoutput}
\begin{verbatim}
id_etudiant | nom     | prenom | date_naissance | filiere
-----------+---------+--------+----------------+--------
1           | Bernard | Alice  | 2004-03-15     | Info
2           | Petit   | Bob    | 2003-11-22     | Info
4           | Leroy   | David  | 2003-01-30     | Info
\end{verbatim}
\end{codeoutput}
\end{exemple}

\section{Projection ($\pi$)}

\begin{definition}[Projection]
La \emph{projection} $\pi_{A_1, A_2, \ldots, A_k}(R)$ retourne la relation
obtenue en ne conservant que les attributs $A_1, \ldots, A_k$ de $R$,
et en \'eliminant les doublons.
\[
\pi_{A_1, \ldots, A_k}(R) = \{ t[A_1, \ldots, A_k] \mid t \in R \}
\]
\end{definition}

\begin{exemple}
Projeter les noms et pr\'enoms des \'etudiants~:
\[
\pi_{\text{nom}, \text{prenom}}(\text{Etudiants})
\]

\begin{codeblock}[title=\textbf{\'Equivalent SQL de la projection}]
\begin{minted}{sql}
SELECT DISTINCT nom, prenom FROM Etudiants;
\end{minted}
\end{codeblock}
\end{exemple}

\begin{attention}[title=\textbf{Projection et doublons}]
En alg\`ebre relationnelle, la projection \'elimine automatiquement les doublons
(car une relation est un ensemble). En SQL, il faut utiliser \texttt{DISTINCT}
explicitement, car SQL travaille avec des \emph{multi-ensembles} (bags).
\end{attention}

\section{Op\'erateurs ensemblistes}

\begin{definition}[Union, Intersection, Diff\'erence]
Soient $R$ et $S$ deux relations de m\^eme sch\'ema (\emph{union-compatibles})~:
\begin{align}
R \cup S &= \{ t \mid t \in R \lor t \in S \} \tag{union}\\
R \cap S &= \{ t \mid t \in R \land t \in S \} \tag{intersection}\\
R - S &= \{ t \mid t \in R \land t \notin S \} \tag{diff\'erence}
\end{align}
\end{definition}

\begin{theoreme}[Propri\'et\'es ensemblistes]
\begin{itemize}
    \item $\cup$ et $\cap$ sont commutatives et associatives.
    \item $-$ n'est ni commutative ni associative.
    \item $R \cap S = R - (R - S)$.
\end{itemize}
\end{theoreme}

\begin{codeblock}[title=\textbf{Op\'erateurs ensemblistes en SQL}]
\begin{minted}{sql}
-- Union (elimine les doublons)
SELECT nom FROM Etudiants WHERE filiere = 'Info'
UNION
SELECT nom FROM Etudiants WHERE filiere = 'Maths';

-- Intersection
SELECT id_etudiant FROM Inscriptions WHERE id_cours = 101
INTERSECT
SELECT id_etudiant FROM Inscriptions WHERE id_cours = 102;

-- Difference
SELECT id_etudiant FROM Inscriptions WHERE id_cours = 101
EXCEPT
SELECT id_etudiant FROM Inscriptions WHERE id_cours = 103;
\end{minted}
\end{codeblock}

\section{Produit cart\'esien ($\times$)}

\begin{definition}[Produit cart\'esien]
Le \emph{produit cart\'esien} $R \times S$ produit toutes les combinaisons possibles
de tuples de $R$ avec des tuples de $S$.
\[
R \times S = \{ (t_r, t_s) \mid t_r \in R \land t_s \in S \}
\]
Si $|R| = n$ et $|S| = m$, alors $|R \times S| = n \times m$.
\end{definition}

\begin{attention}[title=\textbf{Co\^ut du produit cart\'esien}]
Le produit cart\'esien est tr\`es co\^uteux~: il multiplie les cardinalit\'es.
En pratique, on l'utilise rarement seul~; on le combine avec une s\'election
pour obtenir une \emph{jointure}.
\end{attention}

\section{Jointure ($\bowtie$)}

\begin{definition}[Jointure naturelle]
La \emph{jointure naturelle} $R \bowtie S$ combine les tuples de $R$ et $S$
qui ont les m\^emes valeurs sur les attributs communs.
\[
R \bowtie S = \sigma_{\text{attributs communs \'egaux}}(R \times S)
\]
avec \'elimination des colonnes dupliqu\'ees.
\end{definition}

\begin{definition}[Theta-jointure]
La \emph{theta-jointure} $R \bowtie_\theta S$ est d\'efinie par~:
\[
R \bowtie_\theta S = \sigma_\theta(R \times S)
\]
o\`u $\theta$ est une condition quelconque portant sur les attributs de $R$ et $S$.
Cas particulier~: l'\emph{\'equi-jointure} utilise uniquement des \'egalit\'es.
\end{definition}

\begin{exemple}
Trouver les noms des \'etudiants et les intitul\'es des cours auxquels ils sont inscrits~:
\[
\pi_{\text{nom}, \text{intitule}}(\text{Etudiants} \bowtie \text{Inscriptions} \bowtie \text{Cours})
\]

\begin{codeblock}[title=\textbf{Jointure en SQL}]
\begin{minted}{sql}
SELECT e.nom, c.intitule
FROM Etudiants e
JOIN Inscriptions i ON e.id_etudiant = i.id_etudiant
JOIN Cours c ON i.id_cours = c.id_cours;
\end{minted}
\end{codeblock}

\begin{codeoutput}
\begin{verbatim}
nom     | intitule
--------+------------------
Bernard | Bases de donnees
Bernard | Algorithmes
Petit   | Bases de donnees
Petit   | Reseaux
Moreau  | Algorithmes
Moreau  | Statistiques
Leroy   | Bases de donnees
Roux    | Statistiques
\end{verbatim}
\end{codeoutput}
\end{exemple}

\begin{proposition}[Propri\'et\'es de la jointure]
\begin{itemize}
    \item La jointure naturelle est commutative~: $R \bowtie S = S \bowtie R$.
    \item La jointure naturelle est associative~: $(R \bowtie S) \bowtie T = R \bowtie (S \bowtie T)$.
    \item Si $R$ et $S$ n'ont aucun attribut commun, $R \bowtie S = R \times S$.
\end{itemize}
\end{proposition}

\section{Renommage ($\rho$)}

\begin{definition}[Renommage]
L'op\'erateur de \emph{renommage} $\rho_{B/A}(R)$ renomme l'attribut $A$ en $B$
dans la relation $R$. On peut aussi renommer la relation elle-m\^eme~:
$\rho_S(R)$.
\end{definition}

Le renommage est utile pour~:
\begin{itemize}
    \item Rendre deux relations union-compatibles.
    \item R\'ealiser une auto-jointure (jointure d'une relation avec elle-m\^eme).
\end{itemize}

\begin{exemple}
Trouver les paires d'\'etudiants inscrits au m\^eme cours~:
\[
\pi_{i_1.\text{id\_etud}, i_2.\text{id\_etud}}(
    \sigma_{i_1.\text{id\_etud} < i_2.\text{id\_etud}}(
        \rho_{i_1}(\text{Inscriptions}) \bowtie_{\text{id\_cours}} \rho_{i_2}(\text{Inscriptions})
    )
)
\]
\end{exemple}

\section{Division ($\div$)}

\begin{definition}[Division]
Soient $R(A, B)$ et $S(B)$ deux relations. La \emph{division} $R \div S$ retourne
les valeurs de $A$ dans $R$ qui sont associ\'ees \`a \emph{toutes} les valeurs de $B$ dans $S$.
\[
R \div S = \{ t[A] \mid t \in R \land \forall s \in S, (t[A], s[B]) \in R \}
\]
\end{definition}

\begin{theoreme}[Expression de la division]
La division peut s'exprimer en termes des autres op\'erateurs~:
\[
R \div S = \pi_A(R) - \pi_A\big((\pi_A(R) \times S) - R\big)
\]
\end{theoreme}

\begin{exemple}
Trouver les \'etudiants inscrits \`a \emph{tous} les cours du professeur 2~:

Soit $S = \pi_{\text{id\_cours}}(\sigma_{\text{id\_prof}=2}(\text{Cours}))$
(les cours du professeur 2).

$\pi_{\text{id\_etudiant}, \text{id\_cours}}(\text{Inscriptions}) \div S$

donne les \'etudiants inscrits \`a tous ces cours.

\begin{codeblock}[title=\textbf{Division en SQL (NOT EXISTS)}]
\begin{minted}{sql}
SELECT DISTINCT i.id_etudiant
FROM Inscriptions i
WHERE NOT EXISTS (
    SELECT c.id_cours
    FROM Cours c
    WHERE c.id_prof = 2
    AND NOT EXISTS (
        SELECT 1
        FROM Inscriptions i2
        WHERE i2.id_etudiant = i.id_etudiant
        AND i2.id_cours = c.id_cours
    )
);
\end{minted}
\end{codeblock}
\end{exemple}

\section{Arbre alg\'ebrique et optimisation}

Une expression d'alg\`ebre relationnelle peut \^etre repr\'esent\'ee sous forme
d'arbre~: les feuilles sont les relations de base, les n\oe{}uds internes
sont les op\'erateurs.

\begin{theoreme}[R\`egles d'optimisation heuristique]
\begin{enumerate}
    \item \textbf{Descendre les s\'elections}~: appliquer les s\'elections le plus
          t\^ot possible pour r\'eduire la taille des relations interm\'ediaires.
    \item \textbf{Descendre les projections}~: projeter le plus t\^ot possible
          pour r\'eduire le nombre d'attributs.
    \item \textbf{Regrouper s\'election + produit cart\'esien en jointure}~:
          $\sigma_\theta(R \times S) \to R \bowtie_\theta S$.
    \item \textbf{R\'eordonner les jointures}~: commencer par les jointures
          qui produisent les plus petits r\'esultats interm\'ediaires.
\end{enumerate}
\end{theoreme}

\begin{exemple}
Consid\'erons la requ\^ete~: \og noms des \'etudiants en Info inscrits au cours 101 \fg.

Expression na\"ive~:
\[
\pi_{\text{nom}}(\sigma_{\text{filiere='Info'} \land \text{id\_cours}=101}(\text{Etudiants} \times \text{Inscriptions}))
\]

Expression optimis\'ee~:
\[
\pi_{\text{nom}}(\sigma_{\text{filiere='Info'}}(\text{Etudiants}) \bowtie \sigma_{\text{id\_cours}=101}(\text{Inscriptions}))
\]

La version optimis\'ee \'evite le produit cart\'esien complet et filtre au plus t\^ot.
\end{exemple}

\section{R\'esum\'e des correspondances alg\`ebre--SQL}

\begin{formules}[title=\textbf{Correspondance Alg\`ebre Relationnelle $\leftrightarrow$ SQL}]
\begin{center}
\begin{tabular}{ll}
\toprule
\textbf{Alg\`ebre} & \textbf{SQL} \\
\midrule
$\sigma_\theta(R)$ & \texttt{SELECT * FROM R WHERE $\theta$} \\
$\pi_{A_1,\ldots,A_k}(R)$ & \texttt{SELECT DISTINCT $A_1$, \ldots, $A_k$ FROM R} \\
$R \cup S$ & \texttt{R UNION S} \\
$R \cap S$ & \texttt{R INTERSECT S} \\
$R - S$ & \texttt{R EXCEPT S} \\
$R \times S$ & \texttt{SELECT * FROM R, S} \\
$R \bowtie S$ & \texttt{SELECT * FROM R NATURAL JOIN S} \\
$R \bowtie_\theta S$ & \texttt{SELECT * FROM R JOIN S ON $\theta$} \\
$R \div S$ & Double \texttt{NOT EXISTS} \\
$\rho_{B/A}(R)$ & \texttt{SELECT A AS B FROM R} \\
\bottomrule
\end{tabular}
\end{center}
\end{formules}

\section{Int\'egration Python}

\begin{codeblock}[title=\textbf{Simulation de l'alg\`ebre relationnelle avec pandas}]
\begin{minted}{python}
import pandas as pd

# Relations comme DataFrames
etudiants = pd.DataFrame({
    'id_etudiant': [1, 2, 3, 4, 5],
    'nom': ['Bernard', 'Petit', 'Moreau', 'Leroy', 'Roux'],
    'filiere': ['Info', 'Info', 'Maths', 'Info', 'Maths']
})
inscriptions = pd.DataFrame({
    'id_etudiant': [1, 1, 2, 2, 3, 3, 4, 5],
    'id_cours': [101, 102, 101, 103, 102, 104, 101, 104],
    'note': [15.5, 13.0, 12.0, 16.0, 17.5, 14.0, 9.0, 18.0]
})

# Selection : sigma_{filiere='Info'}(Etudiants)
info = etudiants[etudiants['filiere'] == 'Info']
print("Selection :\n", info)

# Projection : pi_{nom}(Etudiants)
noms = etudiants[['nom']].drop_duplicates()
print("\nProjection :\n", noms)

# Jointure naturelle
resultat = pd.merge(etudiants, inscriptions, on='id_etudiant')
print("\nJointure :\n", resultat[['nom', 'id_cours', 'note']])
\end{minted}
\end{codeblock}

\section*{Exercices}

\begin{exercice}
Exprimez en alg\`ebre relationnelle les requ\^etes suivantes, puis traduisez-les en SQL~:
\begin{enumerate}
    \item Les noms des professeurs du d\'epartement ``Informatique''.
    \item Les intitul\'es des cours ayant au moins 4 cr\'edits.
    \item Les noms des \'etudiants inscrits au cours ``Bases de donn\'ees''.
\end{enumerate}
\end{exercice}

\begin{exercice}
Soit $R(A, B, C)$ et $S(C, D, E)$. Simplifiez l'expression~:
\[
\pi_{A, D}(\sigma_{B > 10 \land D = 'x'}(R \bowtie S))
\]
en descendant les s\'elections et projections.
\end{exercice}

\begin{exercice}
Exprimez la division suivante en SQL~: \og trouver les \'etudiants inscrits
\`a tous les cours du d\'epartement Informatique \fg.
\end{exercice}

\begin{exercice}
Dessinez l'arbre alg\'ebrique correspondant \`a l'expression~:
\[
\pi_{\text{nom}}(\sigma_{\text{note} > 15}(\text{Etudiants} \bowtie \text{Inscriptions} \bowtie \text{Cours}))
\]
Puis proposez un arbre optimis\'e en appliquant les r\`egles heuristiques.
\end{exercice}

\begin{exercice}
Prouvez que $R \cap S = R - (R - S)$ en utilisant la d\'efinition ensembliste
des op\'erateurs.
\end{exercice}
